\documentclass[12pt,a4paper]{article}
\usepackage[UTF8]{ctex}
\usepackage{amsmath,amssymb,amsfonts}
\usepackage{graphicx}
\usepackage{booktabs}
\usepackage{hyperref}
\usepackage{float}
\usepackage{geometry}
\usepackage{caption}
\usepackage{subcaption}
\usepackage{enumitem}
\usepackage{listings}
\usepackage{xcolor}
\geometry{left=2.5cm,right=2.5cm,top=2.5cm,bottom=2.5cm}

\title{2019年A题 高压油管的压力控制}
\author{MathFlow Team}
\date{\today}

\begin{document}
\maketitle
\tableofcontents
\newpage


\section{问题重述}


本文研究的是2019年全国大学生数学建模竞赛A题——高压油管的压力控制问题。需要通过控制柱塞泵的泵油频率和脉宽，使油管内压力稳定在100MPa。


\section{模型建立}


油管内压力变化遵循可压缩流体连续性方程：


\begin{equation}
  \frac{dP}{dt} = \frac{(Q_{in} - Q_{out}) \cdot K}{V}
  \label{eq:ode}
\end{equation}


其中K=1400.0MPa为燃油体积模量，V=500cm³为油管容积。


\section{模型求解与结果}


采用遗传算法优化，得到最优泵油频率为62Hz，占空比为0.186。


稳态压力波动为4.92MPa，满足精度要求。

\end{document}
